\chapter{Introduction}
\label{chap:intro}

\section{Motivation} \label{sec:s1}


The strongest support is for familiarity-based personalization---letting users base a memory palace on a space they know---not customization for its own sake.
- Best primary reference: Merriman et al. (2016) found that, in virtual routes through a familiar versus less-familiar campus area, both younger and older adults had better novel-object recognition and egocentric direction memory in the familiar context. No benefit appeared when both areas were equally familiar. This directly supports familiar, user-specific layouts. PubMed record
  Merriman, N. A., Ondřej, J., Roudaia, E., O'Sullivan, C., \& Newell, F. N. (2016). Familiar environments enhance object and spatial memory in both younger and older adults. Experimental Brain Research, 234(6), 1555--1574. https://doi.org/10.1007/s00221-016-4557-0
- Best VR-MoL implementation support: Moll and Sykes state that they chose an apartment because familiarity with an environment's structure had produced slightly better recall (PDF p. 5 / article p. 945). They further report that familiarity with the palace is a significant factor in recall across prior studies (PDF p. 10 / article p. 950). Their participant P7 specifically proposed a simulated version of their own home as more effective (PDF p. 21 / article p. 961). Their study did not experimentally compare personal versus generic layouts, so use it as design rationale rather than proof of the personalization effect. Optimized virtual reality-based Method of Loci memorization.pdf
- Mechanism support: Reggente et al. explain that familiar locations may add personal meaning and emotional context, both plausible routes to stronger recall; their actual experiment found that explicitly binding objects to spatial locations yielded 28\% more recalled objects. This supports making the layout meaningful and spatially stable. The Method of Loci in Virtual Reality\_ Explicit Binding of Objects to Spatial Contexts Enhances Subsequent Memory Recall.pdf
Use this background-ready sentence:
Familiarity-based spatial personalization may strengthen virtual memory palaces by giving learners stable, meaningful retrieval cues. In a virtual-environment study, familiar contexts improved recognition of novel objects and egocentric spatial memory; VR Method-of-Loci research likewise identifies familiarity with the memory palace as an important recall factor.

Important limitation: do not claim that ``customization improves memory'' generally. Your papers support customization that increases familiarity, meaningful cues, and navigable loci. The clean experiment still missing is generic palace versus user-familiar palace versus merely user-modified-but-unfamiliar palace.
Caplan et al. is a useful caution: environments differed, but effects were small and navigation knowledge was not a major determinant of verbal-memory success. So frame layout personalization as a promising, testable contributor---not the sole cause of recall. Article
- Draft the background subsection
- Build the literature table
- Design the experiment

\section{Objective}


\section{Contribution}




\section{Organization}



